Dentro de la unidad de negocio donde se desarrolla el proyecto no existía ningún mecanismo automatizado de apoyo a la creación de materiales: la verificación de duplicados y la asignación de categoría se resolvían mediante inspección manual del gestor sobre el catálogo en SAP, apoyándose en las funciones de búsqueda estándar del módulo de materiales. Estas operan por coincidencia de cadenas y no por similitud, de modo que un material registrado con una abreviatura distinta a la que el gestor consulta permanece invisible, y la eficacia de la verificación termina dependiendo del conocimiento acumulado de cada persona sobre las convenciones históricas del catálogo. No se identificó ningún intento previo, dentro de la organización, de abordar este problema mediante analítica o aprendizaje automático.

Sí existe un antecedente de otra naturaleza. La corporación cuenta con un estándar de nomenclatura para las descripciones de materiales basado en la normativa ISO 8000, anterior a este trabajo, bajo el cual está redactada la totalidad del catálogo histórico analizado. La organización, por lo tanto, ya había intentado gobernar la calidad de su catálogo, pero lo hizo por vía normativa: definiendo una regla y esperando que cada gestor la aplicara de forma consistente en el momento de la creación. Los duplicados y las clasificaciones inadecuadas acumulados a lo largo del tiempo evidencian que una norma que no cuenta con un mecanismo que verifique su cumplimiento en el momento de la captura no basta, por sí sola, para sostener la calidad del dato.

El problema cuenta además con una oferta comercial madura. El propio fabricante del ERP ofrece \textit{SAP Master Data Governance}, un módulo licenciado que incorpora creación de datos maestros basada en flujos de trabajo y detección de duplicados mediante algoritmos de emparejamiento configurables \cite{sapmdg2026}; existe asimismo un mercado de proveedores especializados en catálogos de mantenimiento, reparación y operación \cite{verdantis2026}, con casos documentados de proyectos que han procesado más de un millón de registros \cite{ey2026}. La unidad de negocio no cuenta con ninguno de estos módulos ni servicios contratados. Estas alternativas se descartaron por dos razones. La primera es de modelo de costo: operan bajo esquemas de licenciamiento por usuario o de suscripción a servicio, que implican un desembolso recurrente indefinido, mientras que lo que se buscaba era una solución de pago único cuyo costo se concentrara en el desarrollo y no comprometiera el presupuesto operativo de los años siguientes. La segunda es de ajuste: se trata de plataformas de gobierno integral, concebidas para intervenir el ciclo completo del dato maestro y para operar sobre taxonomías genéricas, cuando lo que se requería era una herramienta a la medida, acotada al punto específico del flujo donde se origina el problema y entrenada sobre el histórico y la clasificación internos de la propia corporación.

En cuanto a los antecedentes técnicos, la clasificación automática de descripciones de producto cuenta con precedentes que se remontan a los primeros sistemas de comercio electrónico entre empresas. \citeA{ding2002} desarrollaron \textit{GoldenBullet}, un sistema de clasificación automática de datos de producto contra la taxonomía UNSPSC —la misma con la que el catálogo de la unidad de negocio mantiene correspondencia—, evaluando modelos de espacio vectorial, vecinos más cercanos y Naive Bayes sobre datos reales de catálogo. Su mejor configuración alcanzó una exactitud del 78\% en la primera predicción y del 88\% considerando las diez primeras sugerencias, discriminando entre 421 categorías, cifras que constituyen un punto de referencia sobre el orden de magnitud esperable en esta clase de problema. Resulta particularmente pertinente la advertencia de los autores en el sentido de que una exactitud en ese rango iguala o supera la calidad de la clasificación humana sobre el mismo material, y que perseguir cifras sensiblemente superiores implicaría sobreajustar el modelo a decisiones humanas que a su vez contienen una tasa de error no despreciable. En el estado actual de la disciplina, \citeA{cheng2024} muestran la vigencia de una arquitectura por etapas, en la que un modelo especializado acota las categorías candidatas y una segunda instancia basada en un modelo de lenguaje resuelve entre ellas.

En el ámbito de la detección de duplicados, \citeA{albayrak2022} pusieron en producción un motor para una plataforma de comercio electrónico que combina métricas de similitud textual con un clasificador entrenado sobre pares de productos etiquetados por expertos, demostrando a gran escala que el enfoque híbrido supera en precisión a los métodos no adaptativos. Su trabajo confirma que la comparación por similitud, y no la coincidencia exacta de cadenas, es el camino viable en catálogos con descripciones de texto libre; señala también que la variante más precisa exige disponer de un conjunto etiquetado de pares duplicado y no duplicado, cuya construcción supone un esfuerzo de anotación experta considerable.

En conjunto, estos antecedentes confirman que tanto la detección de duplicados como la clasificación automática de descripciones de producto cuentan con precedentes consolidados de aplicación exitosa en catálogos a escala industrial, y establecen el rango de desempeño que la literatura reporta como equivalente o superior a la clasificación humana para esta clase de tarea. Lo que no se localizó fue un antecedente que integrara ambas capacidades dentro del flujo operativo de creación de materiales de una organización y que midiera su efecto sobre el tiempo de gestión con datos de producción, que es el espacio en el que se inscribe el presente trabajo.
